\begin{mistake}{2026-08-02}{02}

\begin{problemcontent}
设 \(a>0\)，\(f(t)\) 在 \([0, +\infty)\) 上连续有界，证明：方程 \(\frac{\text{d}x}{\text{d}t} + ax = f(t)\ (t \ge 0)\) 的所有解在 \([0, +\infty)\) 上有界。
\end{problemcontent}

\begin{solutioncontent}
方程为一阶线性非齐次微分方程，由通解公式可得：
\[x(t) = e^{-at} \left( C + \int_0^t f(s) e^{as} \, \text{d}s \right) = C e^{-at} + e^{-at} \int_0^t f(s) e^{as} \, \text{d}s\]
因为 \(f(t)\) 在 \([0, +\infty)\) 上连续有界，存在常数 \(M > 0\) 使得 \(|f(t)| \le M\)。
对解两边取绝对值，利用积分的绝对值不等式进行放缩：
\[\begin{aligned}
|x(t)| &\le |C| e^{-at} + e^{-at} \int_0^t |f(s)| e^{as} \, \text{d}s \\
&\le |C| e^{-at} + e^{-at} \int_0^t M e^{as} \, \text{d}s \\
&= |C| e^{-at} + \frac{M}{a} e^{-at} (e^{at} - 1) \\
&= |C| e^{-at} + \frac{M}{a} (1 - e^{-at})
\end{aligned}\]
由于 \(a > 0\) 且 \(t \ge 0\)，有 \(0 < e^{-at} \le 1\) 及 \(0 \le 1 - e^{-at} < 1\)。
因此：
\[|x(t)| \le |C| + \frac{M}{a}\]
因为对于每个确定的解，\(C\) 为常数，且 \(M, a\) 为正定常数，故所有解在 \([0, +\infty)\) 上均有界。
\end{solutioncontent}

\mistakeknowledge{
\begin{itemize}
  \item \textbf{核心公式}：一阶线性微分方程通解公式 \(x(t) = e^{-\int P(t)\text{d}t} \left( C + \int Q(t) e^{\int P(t)\text{d}t} \text{d}t \right)\)
  \item \textbf{易错点}：定积分表达式中的积分变量需使用虚拟变量（如 \(s\)），绝不可与积分上限 \(t\) 混淆；不可遗漏齐次通解项 \(C e^{-at}\)。
  \item \textbf{关键技巧}：利用定积分绝对值不等式 \(\left| \int_0^t g(s)\text{d}s \right| \le \int_0^t |g(s)|\text{d}s\) 与函数有界性 \(|f(t)| \le M\) 进行放缩。
\end{itemize}}

\end{mistake}
